(i) Under $X \in U(2)$, $\beta \mapsto \bar{X}\beta$ gives $\bar{\beta} \mapsto X\bar{\beta}$, and $\alpha \mapsto X\alpha$ gives $\bar{\alpha} \mapsto \bar{X}\bar{\alpha}$; so the $v$-coefficient $\bar{\beta}$ of $c^{*}$ transforms as the $v$-coefficient $\alpha$ of $c$ and its $w$-coefficient $\bar{\alpha}$ as $\beta$ (the action of Theorem \claimref{synthesis.rung-field-on-the-lift-locus}). A coupling built from $c^{*}$ with the same index contractions as one built from $c$ is therefore invariant whenever that one is. The mass entries: on the sixteen of Proposition \claimref{standard-model.hyp-s-slot-statistics}, by Theorem \claimref{standard-model.yukawa-pattern-two-rung-module} and Lemma \claimref{standard-model.invariant-pairing-five-modes}, the mass form of $a\,v_{4} + b\,w_{4}$ is $b$ times the $10$ form of $w_{4}$ on $(L, e^{c})$, $(Q, d^{c})$ and $a$ times that of $v_{4}$ on $(L, \nu^{c})$, $(Q, u^{c})$; applied to $c$ with $Y$ and to $c^{*}$ with $\tilde{Y}$ this gives $M_{d} = \beta Y + \bar{\alpha}\tilde{Y}$ and $M_{u} = \alpha Y + \bar{\beta}\tilde{Y}$. The check verifies the mass forms of $c_{\lambda}$ and $c_{\lambda}^{*}$ on the module for four frames. The right-handed part: on $(v_{4}, v_{5}, w_{4}, w_{5})$, $[b_{1}, \cdot]$ sends $v_{4} \mapsto -\tfrac12w_{5}$, $v_{5} \mapsto \tfrac12w_{4}$, $w_{4} \mapsto -\tfrac12v_{5}$, $w_{5} \mapsto \tfrac12v_{4}$, and $[b_{2}, \cdot]$ sends $v_{4} \mapsto \tfrac12w_{5}$, $v_{5} \mapsto -\tfrac12w_{4}$, $w_{4} \mapsto -\tfrac12v_{5}$, $w_{5} \mapsto \tfrac12v_{4}$ (computed on the module). So $R_{1}(\alpha, \beta) = \tfrac12(\epsilon\beta, \epsilon\alpha)$ and $R_{2}(\alpha, \beta) = \tfrac{\mathrm{i}}{2}(\epsilon\beta, -\epsilon\alpha)$ with $\epsilon(x_{4}, x_{5}) = (x_{5}, -x_{4})$; conjugating $(\alpha, \beta) \mapsto (\bar{\beta}, \bar{\alpha})$ commutes with both (for $R_{2}$ the factor $\mathrm{i}$ and the sign exchange compensate), $\epsilon\bar{T} = T\epsilon$ for $T \in \su(2)$ gives commutation with the weak action, and $[R_{1}, R_{2}]$ is $\alpha \mapsto \tfrac{\mathrm{i}}{2}\alpha$, $\beta \mapsto -\tfrac{\mathrm{i}}{2}\beta$, the hypercharge. $B(b\psi, \chi) + B(\psi, b\chi) = 0$ for $b = b_{1}, b_{2}$ (computed), so $B(\psi, [b, x]\varphi) = -B(b\psi, x\varphi) - B(\psi, x\,b\varphi)$: the mass form of the rotated rung vector is the mass form rotated on the fermions, and since the conjugation commutes with $R_{1}$, $R_{2}$ the same holds for $x \mapsto B(\cdot, x^{*}\cdot)$; so $Yc + \tilde{Y}c^{*}$ is covariant; $R_{1}$, $R_{2}$ and $Y$ span the commutant of the weak $\su(2)$ of \claimref{synthesis.rung-field-level-needed-off-balance-and-not-right-invariant}(2), and the coupling has the two-matrix form of a complex bidoublet in left-right models (Fact \claimref{synthesis.imported-left-right-bidoublet-two-yukawa-couplings}). The per-sector reading maps $c$ to entries depending on $\beta$ alone, vanishing on the $\alpha$ half and injective on the $\beta$ half; $R_{1}$ carries the $\alpha$ half into the $\beta$ half, so the kernel is not invariant, while the kernel of an equivariant map is: no action of the right-handed $\su(2)$ on the fermions makes the per-sector reading covariant.

(ii) $\tr H_{d} = |\beta|^{2}\|Y\|^{2} + |\alpha|^{2}\|\tilde{Y}\|^{2} + 2\operatorname{Re}(\beta\alpha\tr(Y\tilde{Y}^{\dagger}))$ and $\tr H_{u} = |\alpha|^{2}\|Y\|^{2} + |\beta|^{2}\|\tilde{Y}\|^{2} + 2\operatorname{Re}(\alpha\beta\tr(Y\tilde{Y}^{\dagger}))$. The cross terms are equal, so the difference is $(|\alpha|^{2} - |\beta|^{2})(\|Y\|^{2} - \|\tilde{Y}\|^{2}) = 2m(\|Y\|^{2} - \|\tilde{Y}\|^{2})$, and the sum is $(|\alpha|^{2} + |\beta|^{2})(\|Y\|^{2} + \|\tilde{Y}\|^{2}) + 4\operatorname{Re}(\alpha\beta\tr(Y\tilde{Y}^{\dagger}))$ with $\alpha\beta = Q(c)$ in the rung-$4$ frame.

(iii) With $\tilde{Y} = Y$, $M_{d} = (\beta + \bar{\alpha})Y = zY$ and $M_{u} = (\alpha + \bar{\beta})Y = \bar{z}Y$, so $\bar{z}M_{d} = zM_{u}$ and $H_{u} = H_{d}$, giving equal spectra and $V = 1$. On $c_{\lambda}$ ($\alpha = H^{2}/\lambda$, $\beta = -\lambda$, Theorem \claimref{synthesis.rung-field-on-the-lift-locus}(i)) $z = -\lambda(1 - e^{-2\eta})$, while the per-sector conjugate reading of Corollary \claimref{standard-model.one-vev-vector-relations-breaking-pattern} reads $-\lambda Y$ for both sectors (loc. cit. (5)): they agree up to the relative correction $e^{-2\eta}$.

(iv) Substituting $\alpha = H^{2}/\lambda$, $\beta = -\lambda$ gives the two matrices. At $\lambda = H$, $M_{d} = H(\tilde{Y} - Y) = -M_{u}$, and $m(c_{H}) = 0$ (loc. cit. (4)). For $\lambda \ne H$, put $r = H^{2}/\lambda^{2} = e^{-2\eta}$, so $M_{d} = \lambda(-Y + r\tilde{Y})$ and $M_{u} = \lambda(rY - \tilde{Y})$; then $rM_{d} + r^{2}M_{u} = \lambda(r^{3} - r)Y = -(1 - r^{2})(H^{2}/\lambda)Y$, which is the stated share, and $-\lambda Y = -r^{-1}(H^{2}/\lambda)Y$. For the reading: on the D-term action the vacuum lies on the lift family at $m = m^{*} = g'^{2}\xi/(g^{2} + g'^{2})$, set by the level $\xi$ (Theorem \claimref{synthesis.rung-field-d-term-action-and-tree-masses-on-the-lift-locus}), so (ii) there reads $\tr H_{u} - \tr H_{d} = 2m^{*}(\|Y\|^{2} - \|\tilde{Y}\|^{2})$. Computed order of magnitude, in the proof text only: with $\eta \approx \ln(v/H_{0}) \approx 101.9$ (the registered value with $H = H_{0}$; with $H_{0}\sqrt{\Omega_{\Lambda}}$ it changes by about $0.2$; the factors between $\lambda$ and $v$, $v^{2} = 2N$, are of order one), $e^{2\eta} \approx 10^{88}$, so a top entry carried through $c$ would need a down-type entry of about $10^{88}m_{t}$ from the same coupling, or a cancellation of that size against the coupling to $c^{*}$.

(v) At $c_{\lambda}$ ($\alpha = (H^{2}/\lambda, 0)$, $\beta = (-\lambda, 0)$, $H = 1$ in the check) the direction $\alpha' = (0, a)$, $\beta' = (0, r\bar{a})$ satisfies $\alpha'^{\mathsf T}\beta + \alpha^{\mathsf T}\beta' = 0$, so it is tangent to $Q = -\Lambda/3$; its Hermitian products $\operatorname{Re}(x^{*}\cdot y)$ with the gauge-orbit tangents $(T\alpha, \bar{T}\beta)$, $T \in \{T_{1}, T_{2}, T_{3}, Y\}$, and with the neutral tangent $dc_{\lambda}/d\lambda$ vanish (for $T_{1}$, $T_{2}$ because $\beta_{4}r = -\alpha_{4}$; for $T_{3}$, $Y$ and the neutral line by support). As $a$ ranges over $\C$ this is a real plane in the three-dimensional orthogonal complement of the gauge orbit, orthogonal to the neutral line: the charged plane of Theorem \claimref{synthesis.rung-field-on-the-lift-locus}(iii). The covariant coupling's entries along rung $5$ are $\beta'_{5}Y + \bar{\alpha}'_{5}\tilde{Y} = \bar{a}(rY + \tilde{Y})$ (down) and $\alpha'_{5}Y + \bar{\beta}'_{5}\tilde{Y} = a(Y + r\tilde{Y})$ (up), those of the vev $-\lambda(Y - r\tilde{Y})$ and $\lambda(rY - \tilde{Y})$; $N(\alpha', \beta') = \tfrac12|a|^{2}(1 + r^{2})$ and $N(c_{\lambda}) = \tfrac12\lambda^{2}(1 + r^{2})$, which gives the stated ratios. Under the per-sector reading only $\beta'$ enters, with weight $r^{2}$ relative to the vev. Order of magnitude (from memory, standard): with $Y \sim M_{d}/v$ and $\tilde{Y} \sim M_{u}/v$ the charged pair couples to $t$ and $b$ with strength of order $m_{t}/v$, as in a two-doublet model, so top-decay searches apply.

The check computes the mass forms of $c_{\lambda}$ and $c_{\lambda}^{*}$ on the $32 \times 32$ rational module at four frames, the covariance under a sample Gaussian-rational unitary $X$, the two sum rules and the one-matrix case for three complex $(\alpha, \beta)$ with fixed Gaussian-rational $Y$, $\tilde{Y}$, on the lift family the balanced frame, the share formula at four frames and one instance's traces, the right-handed generators, their commutation relations and the kernel argument, and the charged plane with its coupling ratios at one frame exactly and at four frames against the closed forms.
